\medskip

On considère la suite $\left(u_{n}\right)$ définie par: 
$\left\{
\begin{array}{l}
u_{1}=\dfrac{1}{\e} \\
\text{pour tout entier } n \geqslant 1,\: u_{n+1}=\dfrac{1}{\e}\left(1+\dfrac{1}{n}\right) u_{n}
\end{array}
\right.$

\begin{enumerate}
	\item 
\begin{list}{\textbullet}{}
\item Pour $n=1$: $u_2=u_{1+1}=\dfrac{1}{\e}\left (1+\dfrac{1}{1}\right )u_1 =\dfrac{1}{\e}(2)\dfrac{1}{\e}=\dfrac{2}{\e^{2}}$
\item Pour $n=2$: $u_3=u_{2+1}=\dfrac{1}{\e}\left (1+\dfrac{1}{2}\right )u_2 =\dfrac{1}{\e}\left (\dfrac{3}{2}\right )\dfrac{2}{\e^2}=\dfrac{3}{\e^{3}}$
\end{list}

	\item On considère une fonction écrite en langage Python qui, pour un entier naturel $n$ donné, affiche le terme $u_{n}$. On complète les lignes $L_{2}$ et $L_{4}$ de ce programme.

	\begin{center}
		\begin{tabular}{|c|>{\ttfamily}l|}
			\hline
			$L_{1}$ & def suite(n): \\
			$L_{2}$ & ~~\blue u = 1/e \\
			$L_{3}$ & ~~for i in range(1, n): \\
			$L_{4}$ & ~~~~u = \blue 1/e*(1+1/i)*u\\
			$L_{5}$ & return u \\
			\hline
		\end{tabular}
	\end{center}

	\item On admet que tous les termes de la suite $\left(u_{n}\right)$ sont strictement positifs.
	\begin{enumerate}
		\item 
$n\geqslant 1 \implies \dfrac{1}{n}\leqslant 1 \implies 1+\dfrac{1}{n}\leqslant 2$ donc $1+\dfrac{1}{n}\leqslant \e$

		\item 
$1+\dfrac{1}{n}\leqslant \e$ donc $\dfrac{1}{\e} \left (1+\dfrac{1}{n}\right )\leqslant 1$

Or pour tout $n$, $u_n>0$, donc:
$\dfrac{1}{\e} \left (1+\dfrac{1}{n}\right ) u_n\leqslant u_n$
c'est-à-dire $u_{n+1}\leqslant u_n$.

La suite $\left(u_{n}\right)$ est décroissante.
		
		\item 
Pour tout $n$, $u_n>0$ donc la suite $(u_n)$ est minorée par 0.

On a démontré que cette suite était décroissante donc, d'après le théorème de la convergence monotone, la suite $(u_n)$ est convergente.		
		\end{enumerate}

	\item \begin{enumerate}
		\item Soit $\mathcal{P}_n$ la propriété: $u_{n}=\dfrac{n}{\e^{n}}$.
		
\begin{list}{\textbullet}{}
\item \textbf{Initialisation}

Pour $n=1$: $u_1=\dfrac{1}{\e}$ et $\dfrac{n}{\e^{n}}=\dfrac{1}{\e}$

Donc la propriété est vraie au rang $n=1$.

\item \textbf{Hérédité}

On suppose que la propriété est vraie au rang $n\geqslant 1$; c'est l'hypothèse de récurrence.

$u_{n+1} = \dfrac{1}{\e}\left (1+\dfrac{1}{n}\right ) u_n
= \dfrac{1}{\e}\left (\dfrac{n+1}{n}\right ) \dfrac{n}{\e^n}
= \dfrac{n+1}{\e^{n+1}}$
donc la propriété est vraie au rang $n+1$.
\item \textbf{Conclusion}

La propriété est vraie au rang 1, elle est héréditaire pour tout $n\geqslant 1$, donc, d'après le principe de récurrence, elle est vraie pour tout $n\geqslant 1$.
\end{list}		
		
On a donc démontré que, pour tout $n\geqslant 1$, on a: $u_{n}=\dfrac{n}{\e^{n}}$.

		\item 
On sait que
$\ds\lim _{x\to +\infty} \dfrac{\e^{x}}{x}=+\infty$,
donc		
$\ds\lim _{x\to +\infty} \dfrac{x}{\e^{x}}=0$.

On en déduit que
$\ds\lim _{n\to +\infty} \dfrac{n}{\e^{n}}=0$
et donc que
$\ds\lim _{n\to +\infty} u_n=0$.

	\end{enumerate}
\end{enumerate}

\bigskip
